\section{Step 9: Block-Structure, Order, and Adaptation Analysis}
\label{sec:step9_blocks}

\subsection{Objective}
The objective of Step~9 is to test whether the frontal condition effect observed in the previous
steps is partly explained by the blocked structure of the experiment.

More specifically, this step asks:
\begin{quote}
After explicitly modeling block order, within-block position, and condition order, does the
abstract-versus-concrete effect in the left frontal ROI remain?
\end{quote}

This step is motivated by the fact that the previous inferential models focused on:
\begin{itemize}
    \item question condition;
    \item question length and structure;
    \item ENEM item correctness;
    \item participant and question variability,
\end{itemize}
but did not yet promote the block structure itself to an explicit inferential component.

\subsection{Why This Step Is Needed}
The previous steps showed:
\begin{itemize}
    \item a near-threshold frontal effect in the fixed-window covariate-adjusted model;
    \item a weaker but still positive frontal effect in the duration-aware variable-window model;
    \item repeated null findings in the back ROI;
    \item no robust back-ROI item-pair effects after multiple-comparison correction.
\end{itemize}

Therefore, the next important question is whether the remaining frontal tendency is:
\begin{itemize}
    \item a condition effect;
    \item a block-order effect;
    \item a fatigue or adaptation effect across the session;
    \item an effect of trial position within the block;
    \item or some combination of these.
\end{itemize}

\subsection{Core Analytical Decision}
To keep Step~9 interpretable, the following elements should remain fixed from Step~7:
\begin{enumerate}
    \item the same included participant-session cohort;
    \item the same duration-aware front ROI outcome;
    \item the same long-channel-only rule;
    \item the same front ROI as the primary ROI;
    \item the same short-channel exclusion from inference;
    \item the same question-level covariates:
    \begin{itemize}
        \item log total word count;
        \item ENEM item correctness.
    \end{itemize}
\end{enumerate}

The main change in Step~9 is:
\begin{quote}
explicitly include the block structure of the experiment as fixed effects in the trial-level model.
\end{quote}

\subsection{Design Assumption}
This step assumes the blocked task structure described earlier in the project workflow:
\begin{itemize}
    \item 10 blocks in total;
    \item 3 questions per block;
    \item 5 blocks assigned to one condition and 5 blocks assigned to the other condition.
\end{itemize}

If the implementation file contains a different realized block structure for some participants,
the actual Step~9 trial table should use the realized block metadata rather than the idealized design.

\subsection{Primary Scientific Question}
The primary Step~9 question is:
\begin{quote}
In the left frontal ROI, does the condition effect remain after adjusting for block order,
within-block trial position, condition order, question length, and item difficulty?
\end{quote}

\subsection{Primary Outcome}
The primary outcome should remain the same as in Step~7:
\begin{quote}
the trial-level left frontal ROI HbO response computed from the duration-aware variable window.
\end{quote}

This is retained so that Step~9 isolates the effect of modeling block structure rather than changing
the signal summary again.

\subsection{Primary ROI Definition}
The primary ROI remains the same left frontal ROI used previously.

\subsubsection{Front ROI HbO channels}
\begin{verbatim}
S2_D6 hbo
S2_D4 hbo
S1_D6 hbo
S1_D3 hbo
S5_D6 hbo
S5_D4 hbo
S5_D3 hbo
\end{verbatim}

\subsubsection{Back ROI HbO channels}
The back ROI is retained as a secondary benchmark:
\begin{verbatim}
S4_D2 hbo
S3_D1 hbo
S6_D2 hbo
S6_D1 hbo
S6_D5 hbo
S7_D2 hbo
S7_D5 hbo
S8_D1 hbo
S8_D5 hbo
\end{verbatim}

\subsubsection{Short-separation pairs excluded from inference}
\begin{verbatim}
S1_D8
S2_D9
S3_D10
S4_D11
S5_D12
S6_D13
S7_D14
S8_D15
\end{verbatim}

\subsection{Block Variables to Construct}
Step~9 should create the following block-related variables for each valid trial.

\subsubsection{Global block order}
\begin{equation}
g_{ik} \in \{1,2,\dots,10\}
\end{equation}
This is the order of the block in the full session.

\subsubsection{Within-block trial position}
\begin{equation}
p_{ik} \in \{1,2,3\}
\end{equation}
This is the position of the question inside the block.

\subsubsection{Condition order}
Let:
\begin{equation}
o_i =
\begin{cases}
1 & \text{if the participant saw Abstract blocks first} \\
0 & \text{if the participant saw Concrete blocks first}
\end{cases}
\end{equation}

If the task did not actually vary this order across participants, this variable should be dropped
and documented as non-informative.

\subsubsection{Block half}
Define:
\begin{equation}
h_{ik} =
\begin{cases}
0 & \text{if } g_{ik} \in \{1,2,3,4,5\} \\
1 & \text{if } g_{ik} \in \{6,7,8,9,10\}
\end{cases}
\end{equation}

This variable provides a simple early-versus-late sensitivity analysis.

\subsubsection{Optional cumulative trial order}
For descriptive purposes, it is also useful to define:
\begin{equation}
m_{ik} \in \{1,2,\dots,30\},
\end{equation}
the cumulative trial number within the full session.

This variable is optional and should not replace the primary block variables above.

\subsection{Primary Covariates}
Step~9 should retain the same item-level covariates used in the Step~7 primary model:
\begin{equation}
w_k = z\!\left(\log(\texttt{total\_word\_count}_k)\right)
\end{equation}
and
\begin{equation}
c_k = z(\texttt{correctness}_k).
\end{equation}

\subsection{Primary Condition Variable}
Let:
\begin{equation}
A_{ik} =
\begin{cases}
1 & \text{if the trial is Abstract} \\
0 & \text{if the trial is Concrete}
\end{cases}
\end{equation}

\subsection{Primary Step~9 Model}
Let \(y^{(F)}_{ik}\) denote the Step~7 duration-aware front ROI HbO response for participant \(i\)
on trial \(k\).

The primary Step~9 model should be:
\begin{equation}
y^{(F)}_{ik}
=
\beta_0
+
\beta_1 A_{ik}
+
\beta_2 g_{ik}
+
\beta_3 p_{ik}
+
\beta_4 o_i
+
\beta_5 w_k
+
\beta_6 c_k
+
u_i
+
v_q
+
\varepsilon_{ik},
\end{equation}
where:
\begin{itemize}
    \item \(u_i \sim \mathcal{N}(0,\sigma_u^2)\) is the participant random intercept;
    \item \(v_q \sim \mathcal{N}(0,\sigma_v^2)\) is the question random intercept;
    \item \(\varepsilon_{ik} \sim \mathcal{N}(0,\sigma^2)\) is the residual term.
\end{itemize}

\subsection{Primary Inferential Target}
The primary inferential target remains:
\begin{equation}
\beta_1,
\end{equation}
which represents the adjusted abstract-versus-concrete effect after controlling for:
\begin{itemize}
    \item global block order;
    \item within-block trial position;
    \item condition order;
    \item question length;
    \item question difficulty;
    \item participant differences;
    \item question differences.
\end{itemize}

\subsection{Primary Hypothesis}
The primary hypotheses are:

\paragraph{Null hypothesis}
\begin{equation}
H_0: \beta_1 = 0
\end{equation}

\paragraph{Alternative hypothesis}
\begin{equation}
H_1: \beta_1 \neq 0
\end{equation}

The significance threshold should remain:
\begin{verbatim}
alpha = 0.05
\end{verbatim}

\subsection{Secondary Block-Interaction Models}
After the primary block-adjusted model, Step~9 should fit a small number of secondary interaction models.

\subsubsection{Secondary Model A: Condition $\times$ Global Block Order}
\begin{equation}
y^{(F)}_{ik}
=
\beta_0
+
\beta_1 A_{ik}
+
\beta_2 g_{ik}
+
\beta_3 p_{ik}
+
\beta_4 o_i
+
\beta_5 w_k
+
\beta_6 c_k
+
\beta_7 (A_{ik} \times g_{ik})
+
u_i
+
v_q
+
\varepsilon_{ik}
\end{equation}

This model tests whether the condition effect changes across the session.

\subsubsection{Secondary Model B: Condition $\times$ Within-Block Position}
\begin{equation}
y^{(F)}_{ik}
=
\beta_0
+
\beta_1 A_{ik}
+
\beta_2 g_{ik}
+
\beta_3 p_{ik}
+
\beta_4 o_i
+
\beta_5 w_k
+
\beta_6 c_k
+
\beta_7 (A_{ik} \times p_{ik})
+
u_i
+
v_q
+
\varepsilon_{ik}
\end{equation}

This model tests whether the condition effect changes from the first to the third question within each block.

\subsubsection{Secondary Model C: Condition $\times$ Condition Order}
\begin{equation}
y^{(F)}_{ik}
=
\beta_0
+
\beta_1 A_{ik}
+
\beta_2 g_{ik}
+
\beta_3 p_{ik}
+
\beta_4 o_i
+
\beta_5 w_k
+
\beta_6 c_k
+
\beta_7 (A_{ik} \times o_i)
+
u_i
+
v_q
+
\varepsilon_{ik}
\end{equation}

This model tests whether the condition effect depends on whether a participant encountered
Abstract blocks first or Concrete blocks first.

\subsection{Early-vs-Late Sensitivity Model}
For a simpler interpretive check, Step~9 should also include:
\begin{equation}
y^{(F)}_{ik}
=
\beta_0
+
\beta_1 A_{ik}
+
\beta_2 h_{ik}
+
\beta_3 (A_{ik} \times h_{ik})
+
\beta_4 w_k
+
\beta_5 c_k
+
u_i
+
v_q
+
\varepsilon_{ik}
\end{equation}

This model asks whether the condition effect differs between the first half and second half
of the session.

\subsection{Secondary Back ROI Model}
Because the back ROI has been consistently null in the previous steps, it should remain secondary here.
Still, Step~9 should fit one back-ROI block-adjusted model as a benchmark:
\begin{equation}
y^{(B)}_{ik}
=
\beta_0
+
\beta_1 A_{ik}
+
\beta_2 g_{ik}
+
\beta_3 p_{ik}
+
\beta_4 o_i
+
\beta_5 w_k
+
\beta_6 c_k
+
u_i
+
v_q
+
\varepsilon_{ik}
\end{equation}

\subsection{Fixed-Window Benchmark Sensitivity}
Because Step~6 used a fixed 7--11~s frontal outcome and gave the strongest near-threshold frontal result,
Step~9 should include one benchmark analysis that repeats the primary block-adjusted model
using the Step~6 fixed-window front ROI outcome instead of the Step~7 duration-aware outcome.

This allows a direct comparison between:
\begin{itemize}
    \item duration-aware block-adjusted frontal effects;
    \item fixed-window block-adjusted frontal effects.
\end{itemize}

\subsection{Descriptive Block Diagnostics}
Before fitting the mixed models, Step~9 should include a descriptive block-focused diagnostic section.

At minimum, this should report:
\begin{itemize}
    \item mean frontal ROI response by global block order and condition;
    \item mean frontal ROI response by within-block position and condition;
    \item mean response time by global block order and condition;
    \item number of valid trials by block and condition;
    \item simple plots of early vs late blocks by condition.
\end{itemize}

These diagnostics are important because they make block effects visible before any model is fit.

\subsection{Step-by-Step Workflow}

\subsubsection{Step~9.1: Freeze the Step~7 trial table}
Reuse the same included participant-session cohort and the same duration-aware front ROI outcome from Step~7.

\subsubsection{Step~9.2: Merge block metadata}
For each valid trial, add:
\begin{itemize}
    \item block label;
    \item global block order;
    \item within-block trial position;
    \item participant-level condition order;
    \item block half;
    \item optional cumulative trial index.
\end{itemize}

\subsubsection{Step~9.3: Run descriptive block diagnostics}
Create descriptive summaries and figures showing how the frontal ROI and response time vary across:
\begin{itemize}
    \item global block order;
    \item within-block position;
    \item early vs late session half;
    \item condition order when available.
\end{itemize}

\subsubsection{Step~9.4: Fit the primary block-adjusted front ROI model}
Fit the primary mixed-effects model with:
\begin{itemize}
    \item condition;
    \item global block order;
    \item within-block position;
    \item condition order;
    \item log word count;
    \item ENEM correctness;
    \item random intercept for participant;
    \item random intercept for question.
\end{itemize}

\subsubsection{Step~9.5: Fit the secondary interaction models}
Fit:
\begin{itemize}
    \item condition \(\times\) global block order;
    \item condition \(\times\) within-block position;
    \item condition \(\times\) condition order.
\end{itemize}

\subsubsection{Step~9.6: Fit the early-vs-late sensitivity model}
Fit the simpler block-half model to test whether the condition effect differs between early and late blocks.

\subsubsection{Step~9.7: Fit the secondary back ROI block-adjusted model}
Fit the benchmark back-ROI block-adjusted model.

\subsubsection{Step~9.8: Fit the fixed-window benchmark model}
Repeat the primary block-adjusted model using the Step~6 fixed-window frontal outcome.

\subsubsection{Step~9.9: Summarize model coefficients and conclusions}
Collect:
\begin{itemize}
    \item the adjusted condition coefficient;
    \item the main block coefficients;
    \item the interaction coefficients;
    \item standard errors;
    \item p-values;
    \item confidence intervals;
    \item participant and question random-intercept variances.
\end{itemize}

\subsection{Primary Reporting Quantities}
The primary Step~9 results table must report:
\begin{itemize}
    \item number of included participants;
    \item number of included questions;
    \item number of included trials;
    \item condition coefficient \(\beta_1\);
    \item standard error;
    \item test statistic;
    \item p-value;
    \item 95\% confidence interval;
    \item global block-order coefficient;
    \item within-block position coefficient;
    \item condition-order coefficient;
    \item log word-count coefficient;
    \item ENEM-correctness coefficient;
    \item participant random-intercept variance;
    \item question random-intercept variance.
\end{itemize}

\subsection{Secondary Reporting Quantities}
The secondary Step~9 tables should report:
\begin{itemize}
    \item condition \(\times\) block-order interaction results;
    \item condition \(\times\) within-block-position interaction results;
    \item condition \(\times\) order-condition interaction results;
    \item early-vs-late sensitivity results;
    \item secondary back-ROI model results;
    \item fixed-window benchmark results.
\end{itemize}

\subsection{Primary Conclusion Rule}
The final Step~9 conclusion should be based first on the condition coefficient \(\beta_1\)
from the primary block-adjusted front ROI model.

The interpretation rule should be:
\begin{itemize}
    \item if the adjusted condition coefficient remains significant, conclude that the frontal
    condition effect persists even after explicit control for block structure;
    \item if the adjusted condition coefficient weakens substantially or becomes clearly null,
    conclude that part of the previous frontal tendency may reflect block structure or session order;
    \item if one of the condition-by-block interaction terms is significant, conclude that the
    condition effect varies across the blocked structure of the task and should not be treated
    as block-invariant.
\end{itemize}

\subsection{Interpretation Templates}

\paragraph{Template if the frontal condition effect survives block adjustment.}
In the block-adjusted mixed-effects analysis, the abstract-versus-concrete effect in the left frontal ROI
remained after controlling for global block order, within-block trial position, condition order,
question length, and item difficulty
(\(\beta=\ldots\), \(SE=\ldots\), \(p=\ldots\), 95\% CI \([\ldots,\ldots]\)).
This suggests that the frontal condition effect is not explained primarily by the blocked structure
of the task.

\paragraph{Template if the frontal condition effect weakens after block adjustment.}
In the block-adjusted mixed-effects analysis, the abstract-versus-concrete effect in the left frontal ROI
was reduced and no longer statistically significant after controlling for block structure
(\(\beta=\ldots\), \(SE=\ldots\), \(p=\ldots\), 95\% CI \([\ldots,\ldots]\)).
This suggests that part of the previously observed frontal tendency may be attributable to session order,
within-block position, or other blocked-task effects.

\paragraph{Template if a condition-by-block interaction is significant.}
The block-adjusted analysis identified a significant interaction between condition and block structure
(\(\beta=\ldots\), \(SE=\ldots\), \(p=\ldots\)).
This indicates that the condition effect changes across the blocked progression of the session and
should not be interpreted as constant across time.

\subsection{Expected Outputs}

\subsubsection*{Tables}
\begin{verbatim}
01_step9_trial_block_table.csv
02_step9_block_descriptive_summary.csv
03_step9_primary_front_block_adjusted_model.csv
04_step9_condition_by_blockorder_model.csv
05_step9_condition_by_position_model.csv
06_step9_condition_by_ordercondition_model.csv
07_step9_early_late_sensitivity_model.csv
08_step9_back_block_adjusted_model.csv
09_step9_fixed_window_block_benchmark.csv
10_step9_exclusion_log.csv
11_step9_final_conclusion.csv
\end{verbatim}

\subsubsection*{Figures}
\begin{verbatim}
figures/step9/front_roi_by_global_block_order.png
figures/step9/front_roi_by_withinblock_position.png
figures/step9/response_time_by_global_block_order.png
figures/step9/early_vs_late_condition_plot.png
figures/step9/primary_block_adjusted_condition_effect.png
figures/step9/condition_by_blockorder_interaction_plot.png
figures/step9/condition_by_ordercondition_plot.png
figures/step9/fixed_vs_variable_block_benchmark.png
\end{verbatim}

\subsubsection*{Reports}
\begin{verbatim}
reports/step9/step9_block_analysis_report.tex
reports/step9/tables/primary_block_model_table.tex
reports/step9/tables/interaction_models_table.tex
reports/step9/text/final_conclusion.tex
\end{verbatim}

\subsection{Minimal Figures for the Main Report}
The main Step~9 report should include at least:
\begin{enumerate}
    \item frontal ROI response by global block order, separated by condition;
    \item frontal ROI response by within-block trial position, separated by condition;
    \item response time by global block order, separated by condition;
    \item an early-vs-late condition comparison figure;
    \item a coefficient plot for the primary block-adjusted condition effect;
    \item a benchmark comparison between the duration-aware and fixed-window block-adjusted models.
\end{enumerate}

\subsection{Quality-Control Checks}
Before Step~9 is considered complete, verify that:
\begin{enumerate}
    \item the Step~7 cohort is unchanged unless explicitly re-reviewed;
    \item the trial-level frontal outcome is identical to the Step~7 duration-aware outcome;
    \item block labels and block order are correctly reconstructed;
    \item within-block trial position is valid for all included trials;
    \item condition order is correctly assigned at the participant level, or explicitly documented as unavailable;
    \item the primary model includes both block structure and item covariates;
    \item the primary model converges;
    \item interaction models are clearly labeled as secondary;
    \item the back ROI remains secondary;
    \item the fixed-window benchmark is clearly labeled as a sensitivity analysis.
\end{enumerate}

\subsection{Fallback Rule}
If the primary mixed-effects model fails to converge reliably, the fallback analysis should be:
\begin{quote}
a linear model with participant fixed effects and question fixed effects, keeping the same
front ROI duration-aware outcome and the same block and item covariates.
\end{quote}

\subsection{Minimal Completion Criteria}
Step~9 is complete when:
\begin{itemize}
    \item the trial-level block table has been built;
    \item descriptive block diagnostics have been completed;
    \item the primary block-adjusted front ROI model has been fit;
    \item the secondary block-interaction models have been fit;
    \item the early-vs-late sensitivity has been completed;
    \item the back-ROI benchmark has been fit;
    \item the fixed-window benchmark has been fit;
    \item the primary adjusted p-value has been reported;
    \item a short final conclusion has been generated;
    \item all exclusions and warnings have been documented.
\end{itemize}

\subsection{Short Practical Summary}
In simple terms, Step~9 should:
\begin{quote}
take the Step~7 duration-aware frontal ROI outcome, explicitly add block order and within-block
position into the mixed-effects model, test whether the frontal condition effect survives that adjustment,
and determine whether the remaining signal is better explained by condition, by session order,
or by adaptation across the blocked task structure.
\end{quote}